\documentclass[11pt,a4paper]{article}

\usepackage[T1]{fontenc}
\usepackage[utf8]{inputenc}
\usepackage[margin=2.2cm]{geometry}
\usepackage{booktabs}
\usepackage{longtable}
\usepackage{array}
\usepackage{amsmath}
\usepackage{siunitx}
\usepackage{xcolor}
\usepackage[colorlinks=true,linkcolor=blue!50!black,urlcolor=blue!50!black]{hyperref}
\usepackage{fancyhdr}
\usepackage{enumitem}

% group-digits=integer: without it siunitx also groups the fractional part and
% renders 4.9968e-3 as "4.996,8 x 10^-3".
\sisetup{group-separator={,},group-minimum-digits=4,group-digits=integer}

\definecolor{warn}{RGB}{170,40,20}
\definecolor{ok}{RGB}{20,110,60}
\definecolor{prov}{RGB}{90,90,90}

% provenance tag: every number in this document carries one
\newcommand{\ev}[1]{{\footnotesize\textcolor{prov}{[#1]}}}

% Lightweight callout box. Deliberately built from fcolorbox + minipage rather
% than tcolorbox: tcolorbox loads pgf, whose XeTeX driver is absent from the
% tectonic bundle on this machine.
\newsavebox{\calloutbox}
\newenvironment{callout}[3]% #1 frame colour, #2 background colour, #3 title
{\par\medskip\noindent
 \def\calloutframe{#1}\def\calloutback{#2}%
 \setlength{\fboxrule}{0.9pt}\setlength{\fboxsep}{8pt}%
 \begin{lrbox}{\calloutbox}%
 \begin{minipage}{\dimexpr\textwidth-2\fboxsep-2\fboxrule\relax}%
 \setlength{\parskip}{4pt}%
 \textcolor{#1}{\textbf{#3}}\par\nobreak\smallskip}
{\end{minipage}\end{lrbox}%
 \fcolorbox{\calloutframe}{\calloutback}{\usebox{\calloutbox}}%
 \par\medskip}

\pagestyle{fancy}
\fancyhf{}
\lhead{\footnotesize MAGI --- state of validation}
\rhead{\footnotesize\textcolor{warn}{INTERNAL DRAFT --- NOT FOR SUBMISSION}}
\cfoot{\thepage}
\renewcommand{\headrulewidth}{0.4pt}

\title{\textbf{MAGI: state of validation}\\[2mm]
\large An internal reference for paper writing}
\author{F.\ Monastra (INAF)}
\date{Compiled 14 August 2026}

\begin{document}
\maketitle
\thispagestyle{fancy}

\begin{callout}{warn}{red!4}{Status of this document}
\textbf{This is an internal working reference, not a manuscript and not
submission-ready.} It exists so that numbers entering the proceeding and the
flash talk can be traced to the run that produced them.

Every numeric claim carries a provenance tag \ev{Ex} resolving to
Table~\ref{tab:prov}. Tags point at files and runs on the local machine; they
are \emph{not} literature citations. This document contains no external
references, because none have been verified for it.

Superseded and withdrawn results are retained and marked, not deleted.
Work in progress is marked as such and carries no result.
\end{callout}

\section{Scope}

MAGI is a conditional variational autoencoder that replaces Geant4 transport
from a spacecraft or laboratory mass model down to a closed surface (the
``CryoSphere'') surrounding a cryogenic detector. Instead of tracking primaries
through the full assembly, MAGI generates particle crossings directly on that
surface; those crossings are then injected into the same mass model and
transported the short remaining distance.

The validation question is therefore narrow and well posed: \emph{does a MAGI
crossing produce detector events at the same rate and with the same spectrum as
a real recorded crossing, in the same geometry?}

Three independent sources have been tested: cosmic rays on the SRON X-IFU model
(CR), cosmic rays on the DM1.2 laboratory cryostat, and $^{40}$K
radioactivity on the SRON model.

\subsection{Model version under validation}

All results below are \textbf{v0.8.2}, class
\texttt{CVAE\_MixEnergy\_ContPhi\_TaskAdaptive}: conditional rational-quadratic
spline normalizing-flow continuum with CDF pre-warp, pinned 4\,eV emission
lines, focal-gated line/continuum supervision, and a learnable conditional
coupling prior with zone conditioning. \texttt{latent\_dim}\,=\,8.

v0.8.3 (\texttt{energy\_condition\_geometry=True}, \texttt{latent\_dim}\,=\,12)
was implemented and trained at three seeds and is \textbf{not adopted}
(\S\ref{sec:v083}).

\section{The normalization-free comparison}

The central method, used for every headline number here:

\begin{equation}
R \;=\;
\frac{N^{\mathrm{MAGI}}_{\mathrm{events}} / N^{\mathrm{MAGI}}_{\mathrm{crossings}}}
     {N^{\mathrm{real}}_{\mathrm{events}} / N^{\mathrm{real}}_{\mathrm{crossings}}}
\end{equation}

Both sides are counted \emph{per injected crossing}, through the identical mass
model, binary, physics list and scoring code. $R$ requires no exposure, no
detector area, no flux normalization, and---critically---no crossing efficiency
$\chi$. It is insensitive to every quantity that has historically caused
discordance between MAGI-side and Geant4-side bookkeeping.

The \emph{efficiency factor} quoted in some tables is a different normalization,
retained for continuity with earlier analysis:
\begin{equation}
\varepsilon \;=\;
\frac{N^{\mathrm{MAGI}}_{\mathrm{counts}} / N^{\mathrm{MAGI}}_{\mathrm{shot}}}
     {N^{\mathrm{ref}}_{\mathrm{counts}} / N^{\mathrm{ref}}_{\mathrm{shot}}},
\qquad
R = \varepsilon \times \chi ,
\end{equation}
where MAGI's shot particles are CryoSphere crossings and the reference's are
thrown primaries. A perfect surrogate returns $\varepsilon = 1/\chi$, not 1.

\section{Result 1 --- Cosmic rays, DM1.2 laboratory cryostat}
\label{sec:dm12}

The cleanest measurement available, and the only one with a multi-seed ensemble
at detector level.

\paragraph{Configuration.}
Reference: $2\times10^{8}$ muons thrown isotropically
(\texttt{/gps/ang/type iso}, \texttt{mintheta 0}, \texttt{maxtheta 90}) through
the DM1.2 mass model, yielding \num{999358} CryoSphere crossings and
\num{1733} detector events \ev{E1}. MAGI: \num{1200000} crossings generated
from a 3-seed ensemble (seeds 42, 7, 13; two files per seed) and injected into
the same model, yielding \num{2019} detector events \ev{E2}.
Sphere: centre $(0,0,5)$\,mm, $R=95$\,mm, verified from model metadata at
generation time \ev{E3}.

\begin{table}[h]
\centering
\caption{DM1.2: counts per injected crossing, MAGI vs full simulation.}
\begin{tabular}{lrrr}
\toprule
Band & Reference & MAGI & $R$ \\
\midrule
all deposits            & \num{1733} & \num{2019} & $0.970 \pm 0.032$ \\
\textbf{100--300\,keV (MIP)} & \num{1325} & \num{1589} & $\mathbf{0.999 \pm 0.037}$ \\
300--1000\,keV          & 310  & 323  & $0.868 \pm 0.069$ \\
30--100\,keV            & 42   & 46   & $0.912 \pm 0.195$ \\
0.1--30\,keV (X-IFU)    & 41   & 43   & $0.873 \pm 0.191$ \\
$>$1000\,keV            & 13   & 17   & $1.089 \pm 0.401$ \\
\bottomrule
\end{tabular}
\ev{E2}
\end{table}

\paragraph{Efficiency factor.}
$\chi_{\mathrm{DM1.2}} = \num{999358}/\num{200000000} = \num{4.9968e-3}$, so a
perfect surrogate returns $1/\chi = 200.13$ \ev{E2}.

\begin{table}[h]
\centering
\caption{DM1.2 efficiency factor. The 0.1--30\,keV row is the one directly
comparable with the SRON CR value in \S\ref{sec:cr}.}
\begin{tabular}{lrrr}
\toprule
Band & $\varepsilon$ & ideal $1/\chi$ & $R = \varepsilon\chi$ \\
\midrule
0.1--30\,keV (X-IFU)  & $174.8 \pm 38.2$ & 200.13 & 0.873 \\
1--20\,keV            & $217.9 \pm 56.8$ & 200.13 & 1.089 \\
100--300\,keV (MIP)   & $199.9 \pm 7.4$  & 200.13 & 0.999 \\
all deposits          & $194.2 \pm 6.4$  & 200.13 & 0.970 \\
\bottomrule
\end{tabular}
\ev{E2}
\end{table}

\paragraph{What can and cannot be claimed.}
The MIP band is the spectral test: \num{1325} reference events, $\pm2.7\,\%$,
and $\varepsilon$ within $0.4\,\%$ of its ideal value. The X-IFU band is an
\emph{integrated rate only} --- 41 reference events, $\pm15.6\,\%$ before MAGI
contributes any error --- and must not be presented as a spectral comparison.
More MAGI generation cannot improve it; the reference is the cap.

DM1.2 has a \textbf{single} sensitive detector, so there is no across-detector
spread to quote as a systematic; the $\pm38.2$ is pure Poisson.

\paragraph{Structural checks (all pass).}
Crossing coverage \num{1733}/\num{1733} $= 100.0\,\%$ \ev{E2}; global species
fractions correct to $0.3\,\%$ ($\gamma$ 1.003, $e^{-}$ 1.002, $e^{+}$ 1.039,
$\mu^{-}$ 0.997) \ev{E2}; impact-parameter distribution 0.912--1.010 over
$b<1$ to $b<50$\,mm \ev{E2}.

\paragraph{Known defect: a coherent 4\,\% deposit shift.}
A two-sample Kolmogorov--Smirnov test on all deposits gives $D = 0.0558$,
$p = 0.0057$ \ev{E2}. This is \emph{not} a shape distortion but a uniform
downward shift: deposit quantiles 25--95\,\% run 0.984, 0.960, 0.959, 0.940,
0.962, and the maximum CDF separation falls on the MIP peak at 186\,keV. It
localises to the MIP band ($D=0.0594$, $p=0.012$); the 300--1000\,keV band and
everything below 100\,keV are indistinguishable ($p = 0.81$ and $0.99$).

The cause is path length. Among detector-aimed crossings MAGI's tracks are too
close to normal incidence: median $\sec\theta$ is 0.955 of the reference value,
against a measured deposit-median ratio of 0.960 \ev{E4}. Underlying it is a
\emph{direction-conditional species-mix error}: at $b<10$\,mm MAGI runs
$+8.8\,\%$ muons, $-13.6\,\%$ $e^{-}$, $-7.7\,\%$ $\gamma$, while global
fractions are correct to $0.3\,\%$ \ev{E4}. Per-species energies are correct
(muons 0.99--1.03 at every cut).

\paragraph{Prerequisite finding: a geometry overlap that destroyed 39\,\% of crossings.}
At $R=105$\,mm the CryoSphere intersected \texttt{PlateCu1} (a 150\,mm-radius,
5\,mm-thick copper disc at $z=106.5$, spanning $z\in[104,109]$) by at least
1.29\,mm --- 13$\times$ the shell thickness --- and was placed with overlap
checking disabled. Geant4 navigation inside an overlap is undefined, and
$39.4\,\%$ of detector events had \emph{no recorded crossing}, uniformly
58--63\,\% across five independent runs, $93.5\,\%$ of them primary-muon
deposits \ev{E5}. Moving the sphere to $R=95$\,mm restores $100.0\,\%$
coverage. The SRON sphere was checked and is clean.

Separately, the DM1.2 muon macro had used \texttt{/gps/ang/type focused} ten
times, aiming every muon at the origin (median $\cos$ 0.994 to the origin,
$56.8\,\%$ inside $\cos>0.99$); absolute rates from that configuration were
never physical \ev{E5}. All DM1.2 results here are from the \texttt{iso}
rerun.

\section{Result 2 --- Cosmic rays, SRON X-IFU model}
\label{sec:cr}

\paragraph{Normalization-free.}
$R = \mathbf{0.848 \pm 0.125}$ \ev{E6}, from 51 counts in \num{1600000}
injected real crossings against 500 counts in \num{18500000} injected MAGI
crossings, same geometry and binary. (The uncertainty combines both arms; an
earlier statement of $\pm0.14$ used the reference arm alone.)

\paragraph{Per-detector flux ratio} against the RUN CRYOAC reference, 1--20\,keV:
0.76, 0.90, 0.85, 0.57, 0.81, 0.87 over detectors 1--6, mean \textbf{0.79}
\ev{E7}. Spectral shape is \textbf{indistinguishable} from the reference:
KS $D = 0.048$, $p = 0.20$ \ev{E7}. For contrast, v0.6.2 gave $p = 0.0014$.

\paragraph{Efficiency factor.} Counts in 0.1--30\,keV \ev{E7}:

\begin{table}[h]
\centering
\caption{SRON CR efficiency factor per detector, v0.8.2. Ideal value
$1/\chi_{\mathrm{CR}} = 164.4$.}
\begin{tabular}{lrr}
\toprule
Detector & Reference counts & $\varepsilon$ (v0.8.2) \\
\midrule
1 --- TES array     & \num{28402} & $124.8 \pm 5.6$ \\
2 --- Back CryoAC   & \num{11187} & $142.6 \pm 9.6$ \\
3 --- Lateral       & \num{2922}  & $121.3 \pm 17.3$ \\
4 --- Lateral       & \num{3036}  & $98.1 \pm 15.2$ \\
5 --- Lateral       & \num{2829}  & $122.8 \pm 17.7$ \\
6 --- Lateral       & \num{3039}  & $135.3 \pm 17.9$ \\
\midrule
mean                &             & \textbf{124.1} \\
std across detectors&             & 15.2 \\
\bottomrule
\end{tabular}
\end{table}

Normalized, $R = \varepsilon\chi = \mathbf{0.76}$, consistent with the 0.79 of
the flux ratio and the 0.848 of the normalization-free arm, by three
independent code paths.

\paragraph{The SRON--DM1.2 contrast.} This is the most important open physics
question in the validation:

\begin{center}
\begin{tabular}{lrrr}
\toprule
 & $\varepsilon$ (0.1--30\,keV) & ideal $1/\chi$ & $R$ \\
\midrule
SRON CR, 6 detectors  & $124.1 \pm 15.2$ & 164.4  & \textbf{0.76} \\
DM1.2, single detector& $174.8 \pm 38.2$ & 200.13 & $0.873 \pm 0.191$ \\
DM1.2, MIP band       & $199.9 \pm 7.4$  & 200.13 & $0.999 \pm 0.037$ \\
\bottomrule
\end{tabular}
\end{center}

DM1.2 sits at its ideal value where statistics are good; SRON shows a
$\sim$24\,\% deficit that DM1.2 does not reproduce. The DM1.2 overlap bug was
excluded as the cause --- the SRON sphere does not overlap.

\paragraph{Identified contribution to the SRON deficit.}
An \textbf{inconsistent ingoing-particle cut} \ev{E8}. \texttt{EventAction.cc}
skips outgoing crossings at registration, but the SRON training sets predate
that code: the CR training set contains $10.37\,\%$ outgoing crossings and the
$^{40}$K set $22.53\,\%$, while the DM1.2 set contains $0.00\,\%$. MAGI applies
no ingoing/outgoing selection either, and generates $13.59\,\%$ outgoing against
the training set's $10.37\,\%$ --- accounting for $3.6\,\%$ of the rate deficit.
This also means \textbf{SRON and DM1.2 models are on different conventions.}

\section{Result 3 --- $^{40}$K, SRON BuildingModel}

\paragraph{$R = \mathbf{1.118 \pm 0.373}$, KS $D = 0.2384$, $p = 0.60$}
\ev{E9}: $2\times10^{8}$ injected on both sides into
\texttt{SRON\_CCNCoolerwithXFDM\_noShields\_BuildingModel.gdml}, giving 19 MAGI
counts (\num{9.50e-8} per crossing) against 17 from real recorded crossings
(\num{8.50e-8}). Matched $N$ by construction. \textbf{$^{40}$K passes at
detector level.}

The real-crossing arm recycles \num{1873609} recorded crossings $\sim$107$\times$,
which is sound because Geant4 re-randomizes the interaction on each replay;
continuity was verified from the crossing counts ($\sim$40.15\,M re-registered
per job, 0.80 per injected event, uniform across four jobs) rather than assumed
\ev{E9}.

\paragraph{The 22$\times$ gap is Geant4 bookkeeping, not MAGI.}
Against the 44-days reference the same run reads $0.049 \pm 0.012$. The
reference and the MAGI training set \emph{do not share a $^{40}$K source}: the
reference samples a $2000\times2000\times400$\,mm volume centred at
$z=-1600$\,mm, the training set $750\times750\times50$\,mm at $z=-1250$\,mm ---
a \textbf{57$\times$} smaller sampling volume \ev{E10}. The crossing yield
depends on source geometry and not on a scalar, so the two cannot be chained,
and the existing training set cannot be renormalized into the reference.
MAGI sits with the population it was trained on.

\paragraph{Bookkeeping correction.}
\texttt{NgenPotassiumForMAGI\_v082} was declared \num{580000000}; measured from
the run's own output it is \textbf{\num{538205667}} \ev{E11}. The output file
concatenates two jobs, with exactly one backward \texttt{EventId} jump at row
\num{398400}; estimating each job separately by $N = \max\cdot(n+1)/n$ gives
\num{438145967} and \num{100059700}. The declared value is $7.8\,\%$ high and
outside the $95\,\%$ region of both jobs' order statistics.

The same estimator was run as a control on two other declared counts and
\emph{confirms} them ($^{40}$K v0.7.2 ratio 1.0001; CR v0.8.2 ratio 0.9998),
so it is not systematically low. Every v0.8.2 $^{40}$K \emph{Flower} flux and
efficiency rises by 1.0777; the Detector-1 efficiency goes $3472 \to 3741$.
The detector-level result above is unaffected --- it never uses the constant.

\section{Correction: the ``seed systematic'' is withdrawn}
\label{sec:seed}

\begin{callout}{warn}{red!4}{Withdrawn claim}
Earlier versions of the internal documents reported that MAGI's
detector-relevant energy ratio spans $0.455$--$1.222$ across training seeds, and
recommended quoting a $\pm40\,\%$ single-seed / $\pm15\,\%$ ensembled
systematic. \textbf{That spread is an artefact of the metric and is
withdrawn.} It was never measured on a physical quantity.
\end{callout}

The metric \texttt{energy\_vs\_impact\_parameter} compares median generated to
median real energy among detector-aimed crossings. On the CR source that
statistic estimates the \emph{species mix}, not the energy conditional.

At $b<10$\,mm the real population is \ev{E12}: $\mu^{-}$ $48.008\,\%$ at median
\SI{828.373}{\mega\electronvolt}; $\gamma$ $49.789\,\%$ at median
\SI{0.511}{\mega\electronvolt}; $e^{-}$ $2.092\,\%$ at median
\SI{3.170}{\mega\electronvolt}; aggregate median
\SI{13.288}{\mega\electronvolt}. The two dominant populations are three orders
of magnitude apart, split almost exactly 50/50, and \emph{disjoint} ($\gamma$
99th percentile \SI{30.488}{\mega\electronvolt}, $\mu^{-}$ 1st percentile
\SI{24.072}{\mega\electronvolt}). The 50th percentile of the combined sample
therefore falls in the empty gap between them, and the aggregate median is
close to a step function of the muon fraction.

\paragraph{Demonstration using the real data only.} Reweighting the species mix
of the real $b<10$\,mm crossings, with every energy value left untouched, so
that the energy error is zero by construction \ev{E12}:

\begin{center}
\begin{tabular}{rrrr}
\toprule
$\mu^{-}$ reweight & $\mu^{-}$ fraction & aggregate median & ratio to truth \\
\midrule
0.955 & $46.834\,\%$ & \SI{9.287}{\mega\electronvolt}  & \textbf{0.699} \\
0.980 & $47.686\,\%$ & \SI{11.754}{\mega\electronvolt} & 0.885 \\
1.000 & $47.953\,\%$ & \SI{13.172}{\mega\electronvolt} & 0.991 \\
1.012 & $48.325\,\%$ & \SI{14.982}{\mega\electronvolt} & \textbf{1.128} \\
1.050 & $49.268\,\%$ & \SI{21.316}{\mega\electronvolt} & 1.604 \\
\bottomrule
\end{tabular}
\end{center}

A $\pm5\,\%$ mix shift swings the metric over 0.70--1.60. The measured per-seed
muon reweights at $b<10$\,mm are $1.012 / 1.012 / 0.955$ --- precisely this
range, and precisely the seeds that produced the reported spread.

\paragraph{The energy conditional, resolved per species} \ev{E12}:

\begin{center}
\begin{tabular}{llll}
\toprule
Species & seed 42 & seed 7 & seed 13 \\
\midrule
$\gamma$   & 1.000--1.001 & 1.000--1.001 & 1.000--1.039 \\
$\mu^{-}$  & 0.956--1.003 & 1.043--1.064 & 0.947--0.976 \\
$e^{-}$    & 0.680--1.303 & 0.805--1.005 & 0.555--0.615 \\
\bottomrule
\end{tabular}
\end{center}
(ranges over the cuts all-ingoing, $b<50$, $b<20$, $b<10$\,mm)

Gammas are reproduced \emph{exactly} at every cut and every seed; muons within
$\pm6\,\%$. $e^{-}$ scatters but is 1.9--2.1\,\% of the population with 639 real
events at $b<10$\,mm --- statistics-limited, not a demonstrated defect.

\paragraph{What is real.} A $\pm5\,\%$ seed variation in the species mix
\emph{conditional on direction} ($\mu^{-}$ reweight $1.012/1.012/0.955$ at
$b<10$\,mm, against global fractions correct to $0.3\,\%$) \ev{E12}. This is
genuinely a joint-distribution failure that marginal validation cannot see ---
an order of magnitude smaller than claimed, and in a different channel. The
same signature appears independently on DM1.2 (\S\ref{sec:dm12}).

\paragraph{Consequence for \S\ref{sec:cr}.} The \S4b diagnostic in the
validation note --- ``the model has invented an anti-correlation between energy
and aiming'' --- rests on the same aggregate statistic and is \textbf{not
currently supported}. Its marginals table, aimed-fraction column and counts are
unaffected. The detector-level $R = 0.848 \pm 0.125$ never used the metric and
stands.

\section{v0.8.3: implemented, tested, not adopted}
\label{sec:v083}

\texttt{energy\_condition\_geometry=True} plus \texttt{latent\_dim} $8\to12$,
trained at three seeds \ev{E13}. It did not improve the target metric and
\textbf{regressed the two lowest energy decades consistently across all three
seeds}: 0.001--0.01\,MeV from $0.880/0.921/0.995$ to $0.677/0.715/0.738$;
0.01--0.1\,MeV from $1.686/1.296/1.767$ to $2.216/1.714/1.979$. That regression
is measured on \emph{marginals} and therefore survives the correction in
\S\ref{sec:seed}; it remains a sufficient reason not to adopt v0.8.3.

Since \S\ref{sec:seed} shows the energy conditional was not the broken term,
\texttt{energy\_condition\_geometry} was aimed at a quantity that is not at
fault, which explains why it did not help. The flag remains in
\texttt{core/model.py} defaulting to \texttt{False}, so v0.8.2 is bit-identical
and the experiment is resumable.

\paragraph{Revised priority for the next version.}
\begin{enumerate}[nosep]
  \item Make \texttt{energy\_vs\_impact\_parameter} species-resolved by default
        and refuse an aggregate on multi-species input.
  \item Target the aimed species mix: condition the \emph{type} draw on sampled
        geometry, rather than sampling from a global \texttt{type\_probs}.
  \item Measure the detector-level seed spread directly (in progress,
        \S\ref{sec:inprogress}).
  \item More generated events at small $b$ to settle whether the $e^{-}$
        scatter is real.
\end{enumerate}

\section{Work in progress --- no result yet}
\label{sec:inprogress}

\begin{callout}{blue!40!black}{blue!4}{Running as of 14 August 2026, 08:00 CEST}
\textbf{Three-seed detector-level CR comparison on SRON.} Six jobs, two per
checkpoint (seeds 42, 7, 13), \num{1200000} crossings each --- verified exactly
\num{1200000} records per file with no neutrinos --- injected into
\texttt{SRON\_CCNwithXFDM\_NoShield\_FlowerCryoAC\_fixed.gdml}, the same
geometry as the arm-A reference \ev{E14}.

This is the measurement that has never been taken: every seed claim to date,
withdrawn or not, is crossing-level. Expected sensitivity is $\sim$62 Detector-1
counts per seed ($\pm13\,\%$) and $\sim$112 over all six detectors
($\pm9\,\%$) --- enough to exclude a factor-2 spread decisively, marginal at
$\pm20\,\%$, and \textbf{insufficient} to resolve the $\pm5\,\%$ aimed-mix
effect. ETA $\sim$17\,h.

\textbf{No number from this run may be quoted until it completes.}
\end{callout}

\section{What may and may not be quoted}

\begin{callout}{ok}{green!4}{Quotable}
\begin{itemize}[nosep]
  \item CR, SRON, normalization-free: $R = 0.848 \pm 0.125$.
  \item CR, SRON, spectral shape: KS $p = 0.20$, indistinguishable from RUN CRYOAC.
  \item CR, SRON, efficiency: $124.1 \pm 15.2$ over 6 detectors ($R = 0.76$).
  \item DM1.2, 3-seed ensemble: $R = 0.970 \pm 0.032$ overall,
        $0.999 \pm 0.037$ in the MIP band, $100\,\%$ crossing coverage.
  \item DM1.2 efficiency: $174.8 \pm 38.2$ (X-IFU band), $199.9 \pm 7.4$ (MIP),
        against ideal 200.13.
  \item $^{40}$K, BuildingModel: $R = 1.118 \pm 0.373$, KS $p = 0.60$.
  \item Species fractions reproduced to $0.3\,\%$; impact-parameter
        distribution to within $2$--$9\,\%$.
\end{itemize}
\end{callout}

\begin{callout}{warn}{red!4}{Not quotable}
\begin{itemize}[nosep]
  \item \textbf{Any seed systematic}, in either direction. Not established.
  \item Any aggregate median-E-vs-$b$ number (\S\ref{sec:seed}).
  \item The $^{40}$K 44-days comparison ($0.049$) as a MAGI result --- the
        source geometries differ by 57$\times$ (\S4 above).
  \item DM1.2 X-IFU band as a \emph{spectral} comparison --- 41 reference
        events; it is an integrated rate.
  \item Impact-parameter bins below $b \approx 5$\,mm at present sample sizes.
  \item Absolute rates from any DM1.2 \texttt{focused}-source run.
\end{itemize}
\end{callout}

\section{Open items}

\begin{enumerate}[nosep]
  \item \textbf{The SRON 24\,\% deficit is not fully explained.} $3.6\,\%$ is
        attributed to the ingoing-cut mismatch (\S\ref{sec:cr}); the remainder
        is uncharacterized now that the energy$\times$geometry diagnosis is
        withdrawn.
  \item \textbf{Ingoing-only filter} on both SRON training sets, with $\chi$
        recomputed, then retrain --- makes SRON consistent with DM1.2.
  \item \textbf{Detector-level seed spread} on CR (\S\ref{sec:inprogress}).
  \item \textbf{$e^{-}$ behaviour at small $b$} --- needs generation, not seeds.
  \item \textbf{Speed-up figure} (MAGI generation vs full Geant4) for the talk.
        Measured input: on DM1.2, \texttt{iso} throughput is
        \num{2062} muons/s/core.
  \item \texttt{GeneratedParticlesRecycled} did not fire during $\sim$107$\times$
        recycling --- the warning path in \texttt{PrimaryGeneratorAction.cc} is
        unreachable. Silent recycling is exactly what it exists to catch.
  \item \textbf{DM1.2 fluorescence lines} --- the training set supports
        continuum and geometry but is $\sim$9$\times$ short for lines.
\end{enumerate}

\section{Provenance}

\begin{longtable}{p{1.1cm}>{\raggedright\arraybackslash}p{14.5cm}}
\caption{Evidence tags. All paths are local to the analysis machine.}
\label{tab:prov}\\
\toprule
Tag & Artifact \\
\midrule
\endfirsthead
\toprule
Tag & Artifact \\
\midrule
\endhead
\bottomrule
\endfoot
E1 & \texttt{S1GDML\_DM1.2/allOutputGDML\_DM1\_2\_iso.dat},
     \texttt{allDSCryoSphere\_DM1\_2\_iso.dat} --- the $2\times10^{8}$-muon iso
     reference run. \\
E2 & \texttt{S1GDML\_DM1.2/Analysis/MAGI\_DM1\_2\_iso\_comparison.ipynb},
     executed 13/08/2026 22:57; jobs in
     \texttt{S1GDML\_DM1.2/magi\_iso\_run/job0..5}. Figures:
     \texttt{MAGI\_DM1\_2\_iso\_\{spectrum,crossings,aiming\}.png}. \\
E3 & \texttt{magi\_iso\_run/gen\_s0..5.log} --- the
     ``CryoSphere for reconstruction'' line emitted by
     \texttt{MAGI/scripts/generate\_geant\_source.py}. \\
E4 & Incidence-angle and aimed-species analysis over
     \texttt{/Volumes/X10Pro/tmp/magi\_dm12iso\_s0..5.gntbin} against
     \texttt{allDSCryoSphere\_DM1\_2\_iso.dat}, 13/08/2026. \\
E5 & \texttt{S1GDML\_DM1.2/CRYOSPHERE\_OVERLAP\_FIX.md};
     \texttt{src/GDMLDetectorConstruction.cc} (R 105$\to$95, overlap check on);
     \texttt{gpsMu\_iso.mac}. \\
E6 & \texttt{S1GDMLSRON\_CryoSphereEmission/ab\_run/} arm A vs the v0.8.2
     $1.85\times10^{7}$ run. \\
E7 & \texttt{S1GDMLSRON\_CryoSphereEmission/Analysis/DataAnalysis\_pkg.ipynb},
     sections \emph{RUN CRYOAC vs MAGI} and \emph{Efficiency factor per
     detector}; \texttt{docs/MAGI\_validation\_note.md} \S2--\S3. \\
E8 & \texttt{docs/MAGI\_validation\_note.md} \S5;
     \texttt{src/EventAction.cc} ingoing test applied to the training files. \\
E9 & \texttt{k40\_run/run\_now.sh}, \texttt{k40\_run/run\_real\_bm.sh},
     \texttt{k40\_run/compare\_real\_vs\_magi.txt}. \\
E10& \texttt{Old-Simulations/Sim/S1GDMLSRON/gpsPotassium40.mac} vs
     \texttt{S1GDMLSRONBuilding/gpsPotassium40.mac}. \\
E11& \texttt{athena-build/outputGDML\_K40\_v082.dat} EventId analysis;
     \texttt{DataAnalysis\_pkg.ipynb} cell 4;
     \texttt{MAGI\_validation\_note.md} \S7.6. \\
E12& \texttt{MAGI/tools/species\_resolved\_evsb.py},
     \texttt{MAGI/logs/species\_evsb.log}, 13/08/2026; real reference
     \texttt{MAGI/TrainingData/alloutputDSCryoSphereCR.dat}
     (\num{3285347} crossings, neutrino-free). \\
E13& \texttt{MAGI/tools/train\_v083\_cr\_seed.py},
     \texttt{MAGI/tools/compare\_v082\_v083.py},
     \texttt{MAGI/docs/v0.8.3\_seed\_variance.md}. \\
E14& \texttt{S1GDMLSRON\_CryoSphereEmission/cr\_seed\_run/run\_cr\_seeds.sh},
     launched 14/08/2026 07:54. \textbf{In progress.} \\
\end{longtable}

\vspace{4mm}
\noindent\textbf{Companion documents.}
\begin{itemize}[nosep,leftmargin=*]
  \item \texttt{S1GDMLSRON\_CryoSphereEmission/docs/MAGI\_validation\_note.md}
        --- the authoritative running record, including superseded sections.
  \item \texttt{MAGI/docs/v0.8.3\_seed\_variance.md} --- the withdrawn seed
        systematic and the revised plan.
  \item \texttt{S1GDML\_DM1.2/CRYOSPHERE\_OVERLAP\_FIX.md} --- the overlap
        investigation.
\end{itemize}

\end{document}
